\bookchapter{Algebra essentials}\label{ch:algebra}

\section*{Solving and rearranging equations}
The golden rule: \emph{do the same thing to both sides} of an equation. This lets us solve for an unknown.

\begin{example}
$3x+5=20$. Subtract $5$: $3x=15$. Divide by $3$: $x=5$.
\end{example}

\begin{example}[Rearranging a formula]
From $v=v_0+at$ we solve for $t$: $v-v_0=at$, so $t=\dfrac{v-v_0}{a}$. In the same way, from a neuron equation $\hat y=wx+b$ we get the weight $w=\dfrac{\hat y-b}{x}$.
\end{example}

\section*{Inequalities}
Inequalities are solved like equations, with one exception: \textbf{multiplying or dividing by a negative number flips the inequality sign.} The solution is a whole range of numbers, which we draw on a number line. An open circle means the endpoint is \emph{not} included, a filled circle means it \emph{is} included.

\begin{example}
$-2x+3>7$. Subtract $3$: $-2x>4$. Divide by $-2$ and flip the sign: $x<-2$.
\end{example}

\begin{center}
\begin{tikzpicture}[>=Stealth]
  \draw[->] (-5,0) -- (4,0);
  \foreach \x in {-4,-3,-2,-1,0,1,2,3} \draw (\x,2pt) -- (\x,-2pt) node[below,font=\footnotesize]{$\x$};
  \draw[very thick,blue,->] (-2,0) -- (-4.8,0);
  \draw[thick,fill=white] (-2,0) circle (2.5pt);
  \node[anchor=west] at (4.3,0) {$x<-2$};
  \begin{scope}[yshift=-1.8cm]
    \draw[->] (-5,0) -- (4,0);
    \foreach \x in {-4,-3,-2,-1,0,1,2,3} \draw (\x,2pt) -- (\x,-2pt) node[below,font=\footnotesize]{$\x$};
    \draw[very thick,blue,->] (1,0) -- (3.8,0);
    \fill[blue] (1,0) circle (2.5pt);
    \node[anchor=west] at (4.3,0) {$x\ge1$};
  \end{scope}
\end{tikzpicture}

\captionof{figure}{Solution sets of two inequalities on the number line.}\label{fig:b1_6}
\end{center}

\section*{Systems of two linear equations}
Two equations with two unknowns describe two lines. The solution is the point where the lines cross. We can solve by \emph{substitution} or by \emph{elimination}.

\begin{example}
Solve $y=2x-1$ and $y=-x+5$. Setting the right sides equal: $2x-1=-x+5$, so $3x=6$, $x=2$ and then $y=3$. The lines intersect at $(2,3)$.
\end{example}

\begin{center}
\begin{tikzpicture}
\begin{axis}[gr, width=8cm,height=6.5cm, xmin=-1,xmax=5,ymin=-3,ymax=7, xtick={-1,...,5}, ytick={-3,...,7}, legend pos=south east]
  \addplot[blue,domain=-1:4]{2*x-1};
  \addlegendentry{$y=2x-1$}
  \addplot[red,domain=-1:5]{5-x};
  \addlegendentry{$y=-x+5$}
  \addplot[only marks,mark=*,black,mark size=3pt] coordinates {(2,3)};
  \node[above right,font=\small] at (axis cs:2,3) {$(2,3)$};
\end{axis}
\end{tikzpicture}

\captionof{figure}{The solution of the system is the intersection point of the two lines.}\label{fig:b1_7}
\end{center}

\section*{Quadratic equations}
An equation $ax^2+bx+c=0$ can be solved by factoring or by the \emph{quadratic formula}
\[
x=\frac{-b\pm\sqrt{b^2-4ac}}{2a}.
\]
The expression $D=b^2-4ac$ is the \emph{discriminant}. It tells us how many real solutions exist: $D>0$ two solutions, $D=0$ one solution, $D<0$ none. Geometrically, the solutions are the points where the parabola crosses the $x$-axis.

\begin{center}
\begin{tikzpicture}
\begin{groupplot}[group style={group size=3 by 1, horizontal sep=0.9cm},
  width=4.5cm,height=4.4cm,axis lines=middle,xmin=-2.5,xmax=2.5,ymin=-1.8,ymax=4,
  xtick={-2,2}, ytick={-1,1,3}, grid=both,
  grid style={line width=0.1pt, draw=gray!30}, major grid style={line width=0.2pt, draw=gray!50},
  tick label style={font=\scriptsize}, title style={font=\footnotesize\bfseries},
  every axis plot/.append style={thick,blue}, samples=80]
  \nextgroupplot[title={$D>0$: two roots}]
    \addplot[domain=-2.2:2.2]{x^2-1};
    \addplot[only marks,mark=*,red] coordinates {(-1,0)(1,0)};
  \nextgroupplot[title={$D=0$: one root}]
    \addplot[domain=-2.0:2.0]{x^2};
    \addplot[only marks,mark=*,red] coordinates {(0,0)};
  \nextgroupplot[title={$D<0$: no real root}]
    \addplot[domain=-1.7:1.7]{x^2+1};
\end{groupplot}
\end{tikzpicture}

\captionof{figure}{The discriminant decides how often the parabola touches the $x$-axis.}\label{fig:b1_8}
\end{center}

\begin{example}
$x^2-5x+6=0$ factors as $(x-2)(x-3)=0$, so $x=2$ or $x=3$.
\end{example}

\begin{example}
$2x^2+3x-2=0$ has $D=9+16=25$, so $x=\dfrac{-3\pm5}{4}$, i.e.\ $x=\tfrac12$ or $x=-2$.
\end{example}

\section*{Rules for exponents and logarithms}
\begin{center}
\renewcommand{\arraystretch}{1.7}
\begin{tabular}{ll}
\textbf{Exponents} & \textbf{Logarithms}\\\hline
$a^m\cdot a^n=a^{m+n}$ & $\ln(ab)=\ln a+\ln b$\\
$\dfrac{a^m}{a^n}=a^{m-n}$ & $\ln\dfrac ab=\ln a-\ln b$\\
$(a^m)^n=a^{mn}$ & $\ln(a^k)=k\ln a$\\
$a^0=1,\quad a^{-n}=\dfrac1{a^n}$ & $\ln1=0,\quad\ln e=1$\\
$a^{1/2}=\sqrt a$ & $e^{\ln x}=x,\quad\ln(e^x)=x$
\end{tabular}
\end{center}

\begin{example}
Solve $e^{2x}=7$: take $\ln$ on both sides, $2x=\ln7$, so $x=\frac{\ln7}{2}\approx0.973$.
\end{example}

\begin{aibox}
Training a network often means solving for unknown weights, and the loss function \emph{cross-entropy} uses logarithms. The equation $w_1x_1+w_2x_2+b=0$ describes a line in the plane: the \emph{decision boundary} of a neuron. All of these rely on the rearranging skills above.
\end{aibox}

\section*{Practice problems}
\begin{enumerate}[label=\textbf{\thechapter.\arabic*}]
  \item Solve $5x-7=18$.
  \item Solve $3(x-2)=2x+4$.
  \item Rearrange $v=v_0+at$ for $a$, and $\hat y=wx+b$ for $b$.
  \item Solve the inequality $-2x+3>7$ and draw the solution on a number line.
  \item Solve the system $x+y=7$ and $2x-y=2$. Sketch both lines.
  \item Solve $x^2-5x+6=0$ by factoring.
  \item Solve $2x^2+3x-2=0$ with the quadratic formula. Is it true that $D>0$?
  \item Simplify $\dfrac{2^3\cdot2^4}{2^5}$ and $\ln(e^2\cdot e^3)$.
  \item $\star$ Solve $\ln x+\ln(x-2)=\ln3$. (Careful: which solutions are allowed?)
  \item $\star$ A neuron has weights $w_1=2$, $w_2=-1$ and bias $b=-3$. Show that the points with $2x_1-x_2-3=0$ lie on a line, give its equation in the form $x_2=mx_1+c$, and decide on which side of the line the point $(3,0)$ lies (is $2x_1-x_2-3$ positive or negative there?).
\end{enumerate}

% =====================================================
